\clearpage
\fancyhead[L]{\small VERSCHIEDENE AUFSÄTZE ÜBER ASTRONOMIE}
\SourceLeaf{78}%
\pdfbookmark[1]{Berechnung des Osterfestes}{osterfest1800}
\begin{center}
{\large BERECHNUNG DES OSTERFESTES.\par}
\medskip\hrule\medskip
{\small Monatliche Correspondenz zur Beförderung der Erd- und Himmels-Kunde,\par
herausgegeben vom Freiherrn \textsc{von Zach}. August 1800.\par}
\medskip\hrule
\end{center}
\medskip
Dieser Aufsatz lässt das übliche Verfahren zur Bestimmung des Osterfestes beiseite. Man findet es in jeder Anleitung zur mathematischen Chronologie, und es ist an sich leicht genug, wenn man die Bedeutung und den Gebrauch der üblichen Fachwörter \textit{Goldene Zahl}, \textit{Epakte}, \textit{Ostergrenze}, \textit{Sonnenzirkel} und \textit{Sonntagsbuchstabe} kennt und die nötigen Hilfstafeln vor sich hat. Hier geht es um eine rein analytische Lösung der Aufgabe, die von diesen Hilfsbegriffen unabhängig ist und allein auf den einfachsten Rechenoperationen beruht. Hoffentlich ist die Lösung dem Liebhaber willkommen, der mit jenem Verfahren nicht vertraut ist oder das Osterdatum auf der Stelle bestimmen möchte, wenn ihm die nötigen Hilfsmittel fehlen oder er für das betreffende Jahr keinen Kalender nachschlagen kann, und empfiehlt sich zugleich dem Kenner durch ihre Einfachheit und Geschmeidigkeit. Die folgenden Regeln, die sich jeder, der sie der Mühe wert hält, leicht einprägen kann, gelten für zwei Jahrhunderte, von 1700 bis 1899. Durch geeignete Änderung der darin vorkommenden Konstanten und durch Hinzufügen einer geringfügigen Ausnahme lassen sie sich leicht an jedes andere gegebene Jahrhundert anpassen. Die Ausnahme folgt aus der Einrichtung unseres Kalenders und tritt in diesem Zeitraum zufällig nicht ein.

\SourceLeaf{79}%
\expandafter\def\csname HistoricalNoteText2\endcsname{Der Grund ist, dass $a$ nur 19 verschiedene Werte $(0,\ 1,\ 2\ldots18)$ annehmen kann und folglich auch $d$ nur ebenso viele. Der Wert 29 gehört nicht dazu.}
I. Man teilt die Jahreszahl, für die man Ostern berechnen will, durch 19, durch 4 und durch 7 und nennt die Reste dieser Divisionen der Reihe nach $a$, $b$ und $c$. Geht eine Division auf, setzt man den zugehörigen Rest $=0$. Die Quotienten bleiben ganz unberücksichtigt. Dasselbe gilt für die folgenden Divisionen.

II. Dann teilt man $19a+23$ durch 30 und nennt den Rest $d$.

III. Schließlich teilt man $2b+4c+6d+3$ oder $2b+4c+6d+4$ durch 7, je nachdem, ob das gegebene Jahr zwischen 1700 und 1799 oder zwischen 1800 und 1899 einschließlich liegt, und nennt den Rest $e$.

Dann fällt Ostern auf den $22+d+e^{\mathrm{ten}}$ März oder, wenn $d+e$ größer als 9 ist, auf den $d+e-9^{\mathrm{ten}}$ April.

\textit{Beispiele.} Für das Jahr 1744 erhält man bei der Division von 1744 durch 19 den Rest $15=a$. Die Division durch 4 geht auf, also $b=0$. Bei der Division durch 7 erhält man den Rest $1=c$. Daraus ergibt sich $19a+23=308$, was bei der Division durch 30 den Rest $8=d$ ergibt. Schließlich ergibt $2b+4c+6d+3=55$ bei der Division durch 7 den Rest $6=e$. Ostern fällt also auf den $22+8+6^{\mathrm{ten}}$ März oder den $14-9$, das heißt auf den 5. April.

Für 1800 ergibt sich $a=14$, $b=0$, $c=1$. Dann ist $19a+23=289$, also $d=19$. Weiter ist $2b+4c+6d+4=122$, also $e=3$. Somit fällt Ostern auf den $19+3-9$, das heißt auf den 13. April.

Für 1818 ist $a=13$, $b=2$, $c=5$. Dann ist $19a+23=270$, also $d=0$. Weiter ist $2b+4c+6d+4=28$, also $e=0$. Somit fällt Ostern auf den 22. März.

Im letzten Beispiel fällt Ostern auf den frühestmöglichen Tag, denn $d$ und $e$ haben hier offensichtlich ihre kleinstmöglichen Werte. Andererseits ist ersichtlich, dass Ostern nie später als am $22+29+6^{\mathrm{ten}}$ März, das heißt am 26. April, eintreten kann, da $d$ höchstens 29 und $e$ höchstens 6 werden kann. Im achtzehnten und neunzehnten Jahrhundert kann aber niemals $d=29$ werden\HistoricalNote{2}. Der späteste Ostertermin in diesem Zeitraum ist daher der 25. April. Auf ihn fällt Ostern, wenn zugleich $d=28$ und $e=6$ gilt. Diese beiden Bedingungen treffen in den Jahren 1734 und 1886 zusammen. In anderen Jahrhunderten könnte zwar $d=29$ werden. %
\SourceLeaf{80}%
Doch gerade dann greift die oben erwähnte Ausnahme, durch die der Wert von $d$ wieder auf 28 herabgesetzt wird, sodass der 25. April der spätestmögliche Ostertermin überhaupt ist. Dies weiter auszuführen würde hier zu weit führen.

Die Analyse, mit der man die obige Formel findet, beruht eigentlich auf Prinzipien der \textit{höheren Arithmetik}, für die ich gegenwärtig noch keine Schrift anführen kann. Daher lässt sie sich hier freilich nicht in ihrer ganzen Einfachheit darstellen. Das Folgende wird jedoch ausreichen, um sich eine Vorstellung von der Grundlage der Regeln zu machen und sich von ihrer Richtigkeit zu überzeugen.

I. Die Goldene Zahl eines Jahres unserer Zeitrechnung ist bekanntlich der Rest, den man erhält, wenn man zur Jahreszahl 1 addiert und die Summe durch 19 teilt. Geht die Division auf, muss der Rest jedoch $=19$ gesetzt werden. Daraus folgt leicht, dass $a+1$ die Goldene Zahl des gegebenen Jahres ist.

II. Die Ostergrenze, das heißt der Tag des Ostervollmonds, fällt im 18. und 19. Jahrhundert für ein Jahr mit der Goldenen Zahl 1 auf den 13. April. Dann fällt sie im ganzen Zyklus von 19 Jahren, also bis einschließlich des Jahres mit der Goldenen Zahl 19, in jedem Jahr \textit{entweder} 11 Tage früher \textit{oder} 19 Tage später als im unmittelbar vorangehenden Jahr, je nachdem, ob sie in diesem \textit{entweder} in den April \textit{oder} in den März gefallen war, wovon man sich leicht anhand einer Tafel der Ostergrenzen überzeugen kann. Im Jahr mit der Goldenen Zahl 2 fällt sie also auf den 2. April, im darauffolgenden auf den 22. März, im Jahr mit der Goldenen Zahl 4 auf den 10. April und so weiter. Daraus folgt, dass die Ostergrenze nie \textit{vor} dem 21. März und nie \textit{nach} dem 19. April fällt. Nimmt man also an, sie falle in dem Jahr mit der Goldenen Zahl $a+1$ auf den $21+D^{\mathrm{ten}}$ März (wobei man die Tage des Aprils als Fortsetzung des März zählt), so liegt $D$ stets zwischen 0 und 29 einschließlich. Für $a=0$ ist also $D=23$, für $a=1$ ergibt sich $D=23-11$, für $a=2$ ergibt sich $D=23-2\times11$, für $a=3$ ergibt sich $D=23-2\times11+19$ und so weiter. Allgemein gilt $D=23-11p+19q$, wobei $p$ und $q$ durch die Bedingungen bestimmt werden, dass $p+q=a$ gilt und $D$ zwischen 0 und 29 einschließlich liegt. Folglich ergibt sich $D=23+19a-30p$, woraus man leicht schließt, dass $D$ der Rest ist, den man erhält, wenn man $23+19a$ durch 30 teilt. Also gilt $D=d$, oder die Ostergrenze fällt auf den $21+d^{\mathrm{ten}}$ März.

\SourceLeaf{81}%
III. Ostern selbst fällt nur auf den \textit{ersten} Sonntag \textit{nach} der Ostergrenze, also mindestens einen und höchstens sieben Tage später, mithin gewiss nicht vor dem $22+d^{\mathrm{ten}}$ März. Nimmt man also an, Ostern falle auf den $22+d+E^{\mathrm{ten}}$ März, so liegt $E$ zwischen 0 und 6 einschließlich und muss durch die Bedingung bestimmt werden, dass dieser Tag ein Sonntag ist. Diese Bedingung lässt sich rein arithmetisch ausdrücken. Der Zeitraum zwischen dem $22+d+E^{\mathrm{ten}}$ März des gegebenen Jahres und irgendeinem festgelegten Sonntag muss eine durch 7 teilbare Anzahl von Tagen (eine ganze Anzahl von Wochen) umfassen. Man muss also einen bestimmten Sonntag festlegen. Ich wähle dazu den 21. März 1700. Nennt man nun die Zahl des gegebenen Jahres $A$ und bezeichnet mit $i$ die Anzahl der Schaltjahre zwischen 1700 und dem Jahr $A$, dieses eingeschlossen, wenn es ein Schaltjahr ist, so ist $i$ zugleich die Anzahl der Schalttage, die zwischen dem 21. März 1700 und Ostern des Jahres $A$ liegen, und die Anzahl \textit{aller} Tage vom 21. März 1700 bis zum $22+d+E^{\mathrm{ten}}$ März des Jahres $A$
\[
=1+d+E+i+365(A-1700)
\]
Ebenso leicht ist zu sehen, dass zwischen 1700 und 1799 gilt
\[
i=\tfrac14(A-b-1700)
\]
zwischen 1800 und 1899 dagegen
\[
i=\tfrac14(A-b-1700)-1
\]
Zur Bestimmung von $E$ hat man also die Bedingung, dass
\[
1+d+E+365(A-1700)+\tfrac14(A-b-1700)
\]
oder
\[
d+E+365(A-1700)+\tfrac14(A-b-1700)
\]
durch 7 teilbar sein muss, je nachdem, ob das Jahr zwischen 1700 und 1799 oder zwischen 1800 und 1899 liegt. Auch wenn man ein Vielfaches von 7 zu dieser Zahl addiert oder davon abzieht oder wenn man die Zahl von einem Vielfachen von 7 abzieht, muss also eine durch 7 teilbare Zahl herauskommen. Um den Bruch zu beseitigen, addiere ich zunächst $\frac74(A-b-1700)$, was, wie man leicht sieht, durch 7 teilbar ist. Daraus erhalte ich
